\caption{\textbf{Task, recordings and task-aligned feature vectors of $\sim$50,000 neurons.}
\textbf{a} The IBL task and recordings. A mouse keeps the wheel still (quiescence) to start a trial; a visual grating appears left or right, together with an auditory cue, and the mouse turns the wheel to move it to the centre; correct choices are rewarded with water (feedbk1), errors with white noise and a timeout (feedbk0). Neuropixels insertions from 139 mice densely tile the left hemisphere (54,719 neurons passing quality control). Timeline: the peri-event windows (PETH types) that make up each neuron's feature vector, aligned to stimulus onset, movement onset and feedback, plus rest and quiescence.
\textbf{b} Grand average over all neurons ($z$-scored responses, all trials) of each PETH type on a real time axis; movement and feedback windows at the session-averaged event times (movement 253\,ms after the stimulus, feedback 451\,ms after movement; 345 sessions); rest shown next to quiescence (axis break).
\textbf{c} The grand-average feature vector: the 21 PETH types concatenated, as used for every neuron.
\textbf{d} Schematic: how homogeneous is task-related response diversity? Left, every region specialised for one response type (one colour per region); right, every region a mixture, here with the measured response-type proportions (25 k-means clusters, Figs.~2--3) of 12 densely recorded regions. Insets, two example response types (cluster-mean feature vectors).
\textbf{e} Brain-wide latency tiling: everywhere? along the hierarchy? All neurons with at least one spike in the 150\,ms after a left stimulus ($\mathrm{L_sL_cL_b,s}$; 52,975 of 54,719), sorted by the time of their peak (left; each row scaled to [0,~1]) and six example neurons with peaks at evenly spaced latencies (right). Sorting by each neuron's own peak yields a diagonal even for noise; the cross-validated test is in Fig.~5.}
\label{fig:task_data}
